\documentclass[11pt]{article}

\usepackage[margin=1.05in]{geometry}
\usepackage{amsmath,amssymb}
\usepackage{booktabs}
\usepackage{enumitem}
\usepackage{xcolor}
\usepackage[ruled,vlined,linesnumbered]{algorithm2e}
\usepackage[colorlinks=true,linkcolor=blue!45!black,urlcolor=blue!45!black]{hyperref}

\newcommand{\R}{\mathbb{R}}
\newcommand{\E}{\mathbb{E}}
\newcommand{\Var}{\operatorname{Var}}
\newcommand{\tr}{\operatorname{tr}}
\newcommand{\xmap}{x_{\text{\tiny MAP}}}
\newcommand{\code}[1]{\texttt{\small #1}}

\title{\bfseries Whitened Unadjusted Langevin\\[2pt]
\large A sampler for high-dimensional MAGI posteriors}
\author{}
\date{\today}

\begin{document}
\maketitle

\begin{abstract}
\noindent
This note documents \code{MSVGD.whitened\_ula}, an alternative to the SVGD solver for
high-dimensional ODE-inference posteriors. The method reparameterizes the target by the inverse
square root of the Hessian at the mode and runs Langevin dynamics in those coordinates,
Metropolis-adjusted by default. Two
properties make the combination work where SVGD does not: the reparameterization removes the
anisotropy that defeats any scalar-parameterized sampler, and Langevin's Brownian term supplies
the entropy that SVGD's kernel repulsion fails to supply in high dimension. A side effect of
whitening is that the Langevin discretization bias becomes isotropic and known in closed form,
so it can be divided out exactly. On a $d=325$ FitzHugh--Nagumo posterior the method is
statistically indistinguishable from an 8-chain, 64{,}000-draw NUTS reference at roughly
$1/35$ of the cost. We characterize the regime in which it fails---weakly identified posteriors, predictable from
the Hessian spectrum before any sampling---and show that neither the adjusted nor the unadjusted
variant is safe there, for opposite reasons: unadjusted Langevin is transient on the
super-quadratic potentials that polynomial ODE vector fields produce, while the Metropolis
correction cures the transience but is not geometrically ergodic on the same tails, so chains
started outside the bulk freeze instead. Both failures are detected and raised.
\end{abstract}

\section{Motivation}

Stein Variational Gradient Descent evolves an ensemble $\{x_i\}_{i=1}^k$ by a deterministic flow
combining a drift toward high density with a kernel repulsion intended to reproduce the target's
spread. On the posterior considered here it collapses: the ensemble converges to a distribution
far narrower than the target, and does so in a specific, directional way.

The collapse is not a matter of tuning. Initialized at an ensemble that is already correct---exact
draws from a Gaussian statistically indistinguishable from the reference---standard SVGD drives
the Stein-identity dispersion diagnostic
\begin{equation}
  R \;=\; -\frac{1}{k\,d}\sum_{i=1}^{k}\,(x_i-\bar{x})\cdot\nabla\log p(x_i)
  \;\longrightarrow\; 1 \quad\text{under the target}
  \label{eq:steinR}
\end{equation}
from $1.15$ down to $0.03$. Smaller step sizes only slow the decay. The fixed point of the flow
is itself wrong, so no initialization, bandwidth rule, or annealing schedule recovers it.

\section{The geometry of the target}

Write $U(x) = -\log p(x)$ for the negative log posterior on $\R^d$, let $\xmap$ be its minimizer,
and let
\begin{equation}
  H \;=\; \nabla^2 U(\xmap) \;=\; -\nabla^2 \log p(\xmap).
\end{equation}
For the $d=325$ FitzHugh--Nagumo/MAGI posterior, $H$ is positive definite with eigenvalues
spanning
\[
  \lambda_{\min} = 5.31, \qquad \lambda_{\max} = 9.97\times10^{4},
  \qquad \kappa = \lambda_{\max}/\lambda_{\min} = 1.88\times 10^{4}.
\]
Geometrically the posterior is a long, thin ellipsoid: broad in the low-frequency components of
the latent trajectory, extremely tight in the high-frequency components that the Gaussian-process
prior pins down.

This is fatal for any sampler whose behaviour is governed by a single scalar---one step size, or
one kernel bandwidth. A value appropriate to the tight directions is far too small for the broad
ones, and conversely. Section~\ref{sec:analysis} makes the cost precise: it is a factor of
$\kappa$.

\section{The Laplace metric}

If the posterior were exactly Gaussian its covariance would be $H^{-1}$. Take the symmetric
square root $L = H^{-1/2}$ and reparameterize
\begin{equation}
  x \;=\; \xmap + L\,y, \qquad L = L^{\!\top} = H^{-1/2}.
  \label{eq:reparam}
\end{equation}
Under the quadratic model $\log p(x) = -\tfrac12 (x-\xmap)^{\!\top} H (x-\xmap) + c$, substituting
\eqref{eq:reparam} gives
\begin{equation}
  \log \tilde p(y) \;=\; -\tfrac12\, y^{\!\top} H^{-1/2} H H^{-1/2} y + c
  \;=\; -\tfrac12\, y^{\!\top} y + c,
\end{equation}
so the whitened target is $\mathcal{N}(0, I_d)$: a unit sphere. All of the anisotropy has been
absorbed into the change of variables, and a single scalar step size is now appropriate in every
direction. For a target that is only approximately Gaussian, $\tilde p$ is correspondingly only
approximately isotropic, which is all that is required.

\section{Langevin dynamics in whitened coordinates}

The overdamped Langevin diffusion
$\mathrm{d}y_t = \nabla \log \tilde p(y_t)\,\mathrm{d}t + \sqrt{2}\,\mathrm{d}W_t$
has $\tilde p$ as its stationary distribution. Discretized with step size $\varepsilon$ this is
the unadjusted Langevin algorithm (ULA),
\begin{equation}
  y_{t+1} \;=\; y_t + \varepsilon\,\nabla_y \log \tilde p(y_t) + \sqrt{2\varepsilon}\, z_t,
  \qquad z_t \sim \mathcal{N}(0, I_d).
  \label{eq:ula}
\end{equation}

Two remarks.

\paragraph{The chain rule is free.} Because $L$ is symmetric,
\begin{equation}
  \nabla_y \log \tilde p(y) \;=\; \Big(\tfrac{\partial x}{\partial y}\Big)^{\!\top}
  \nabla_x \log p(x) \;=\; L^{\!\top}\nabla_x \log p(x) \;=\; L\,\nabla_x \log p(x).
\end{equation}
The implementation therefore reuses the existing vectorized gradient of the \emph{original}
target and post-multiplies by $L$. No transformed log-density and no second automatic
differentiation pass are needed. (The Jacobian of \eqref{eq:reparam} is constant, so it
contributes nothing to the gradient.)

\subsection{The Metropolis correction}
\label{sec:mala}

\eqref{eq:ula} is biased at finite $\varepsilon$, and Section~\ref{sec:nongauss} shows it is
also unstable on the potentials ODE models produce. Accepting the proposal
$y' = y + \varepsilon\,\nabla\log\tilde p(y) + \sqrt{2\varepsilon}\,z$ with probability
$\min(1,e^{\Lambda})$, where
\begin{equation}
  \Lambda \;=\; \log\tilde p(y') - \log\tilde p(y) + \tfrac12\|z\|^2
  - \frac{\big\|\,y'-y+\varepsilon\nabla\log\tilde p(y')\,\big\|^2}{4\varepsilon},
  \label{eq:mala}
\end{equation}
gives MALA, which is exact for any $\varepsilon$. The forward proposal term collapses to
$\tfrac12\|z\|^2$ because the noise is already in hand, so only the reverse term must be
assembled; carrying $(\log\tilde p,\nabla\log\tilde p)$ between steps costs one
\emph{value-and-gradient} evaluation per step, the same as plain ULA. Non-finite proposals are
rejected automatically, since any comparison against a NaN ratio is false. This is the default.

Because MALA is exact at any step size, $\varepsilon$ affects only efficiency, and it can be
tuned during a discarded warmup by Robbins--Monro on
$\log\varepsilon \leftarrow \log\varepsilon + \gamma(\hat a - a^\star)$ with
$a^\star = 0.574$, the high-dimensional MALA optimum \cite{robertsrosenthal}. A constant gain
suffices because $\hat a$ is averaged over all $k$ chains and so is already precise each step.
In practice the adaptation reaches the same $\varepsilon$ from initial values differing by an
order of magnitude, which removes the step size as a tuning decision.

\paragraph{The noise is the entropy.} This is the conceptual difference from SVGD. SVGD is
deterministic and must manufacture the target's spread through a kernel repulsion term;
Section~\ref{sec:gram} shows that this term silently vanishes in high dimension. ULA instead
carries an explicit Brownian term whose stationary distribution is the target by construction.
There is no kernel and no bandwidth, and the $k$ ``particles'' are $k$ \emph{independent} chains,
so the sampler is embarrassingly parallel.

\section{Analysis}
\label{sec:analysis}

Both of the following results follow from the exactly solvable case: a single eigendirection of a
Gaussian target with precision $\lambda$, i.e.\ variance $1/\lambda$. There
$\nabla \log p(y) = -\lambda y$ and \eqref{eq:ula} becomes the AR(1) recursion
\begin{equation}
  y_{t+1} \;=\; (1-\varepsilon\lambda)\, y_t + \sqrt{2\varepsilon}\, z_t.
  \label{eq:ar1}
\end{equation}

\subsection{Stability and mixing: why whitening is necessary}

Recursion \eqref{eq:ar1} is stable iff $|1-\varepsilon\lambda| < 1$, i.e.
\begin{equation}
  0 < \varepsilon < 2/\lambda .
\end{equation}
This must hold simultaneously in every direction, so the step size is capped by the
\emph{stiffest}: $\varepsilon < 2/\lambda_{\max}$. Meanwhile the autocorrelation in a direction of
precision $\lambda$ decays as $(1-\varepsilon\lambda)^t$, giving a relaxation time of roughly
$\tau_\lambda \approx 1/(\varepsilon\lambda)$ steps, which is worst in the \emph{softest}
direction. Writing $\varepsilon = c/\lambda_{\max}$ with $c<2$,
\begin{equation}
  \tau_{\max} \;\approx\; \frac{1}{\varepsilon \lambda_{\min}}
  \;=\; \frac{\lambda_{\max}}{c\,\lambda_{\min}} \;=\; \frac{\kappa}{c}.
\end{equation}
An unwhitened sampler therefore needs a number of steps proportional to the condition number
$\kappa$. At $\kappa = 1.88\times 10^4$ the stable step size is $\varepsilon < 2\times 10^{-5}$ and
the soft directions need $\mathcal{O}(10^4)$ steps per relaxation time.

After whitening every $\lambda$ equals $1$, so $\varepsilon$ may be $\mathcal{O}(10^{-2})$ and
$\tau = 1/\varepsilon \approx 100$ steps \emph{in all $d$ directions at once}. The dependence on
$\kappa$ is removed entirely. This is why the whole run costs seconds.

\subsection{Discretization bias: exact and removable}

ULA is biased at any finite step size. Taking variances in \eqref{eq:ar1}, the stationary variance
$s_\lambda$ satisfies $s_\lambda = (1-\varepsilon\lambda)^2 s_\lambda + 2\varepsilon$, hence
$s_\lambda\,\varepsilon\lambda\,(2-\varepsilon\lambda) = 2\varepsilon$ and
\begin{equation}
  s_\lambda \;=\; \frac{1}{\lambda}\cdot\frac{1}{1 - \varepsilon\lambda/2},
  \qquad\text{so}\qquad
  \frac{\Var_{\text{ULA}}}{\Var_{\text{target}}} \;=\; \frac{1}{1-\varepsilon\lambda/2}.
  \label{eq:bias}
\end{equation}
Untransformed, this inflation is \emph{anisotropic}: it varies across directions with $\lambda$,
and diverges as $\varepsilon\lambda \to 2$. It is therefore not correctable without knowing the
whole spectrum.

After whitening, $\lambda \equiv 1$ and \eqref{eq:bias} collapses to a single scalar that is the
same in every direction,
\begin{equation}
  \boxed{\;\frac{\Var_{\text{ULA}}}{\Var_{\text{target}}} \;=\; \frac{1}{1-\varepsilon/2}\;}
  \label{eq:isobias}
\end{equation}
which contains no unknowns. Under MALA there is no such bias to remove; for the unadjusted
variant it can be removed exactly by rescaling the whitened deviations,
\begin{equation}
  y_i \;\longleftarrow\; \bar{y} + (y_i - \bar{y})\sqrt{1-\varepsilon/2}\,,
  \label{eq:correction}
\end{equation}
a correction that is exact in the Gaussian limit the method already assumes and that introduces
no tuning constant. Table~\ref{tab:bias} verifies \eqref{eq:isobias} on a $40$-dimensional
anisotropic Gaussian, and Table~\ref{tab:corr} shows the effect of \eqref{eq:correction}.

\begin{table}[t]
\centering
\begin{tabular}{rrrr}
\toprule
$\varepsilon$ & measured $R$ & $1/(1-\varepsilon/2)$ & rel.\ error \\
\midrule
0.200 & 1.1087 & 1.1111 & 0.0022 \\
0.100 & 1.0495 & 1.0526 & 0.0030 \\
0.050 & 1.0236 & 1.0256 & 0.0020 \\
0.020 & 1.0086 & 1.0101 & 0.0015 \\
0.010 & 1.0042 & 1.0050 & 0.0008 \\
\bottomrule
\end{tabular}
\caption{The predicted isotropic inflation \eqref{eq:isobias} against the measured Stein
diagnostic \eqref{eq:steinR}. Target: $d=40$ Gaussian with $\operatorname{cond}(\Sigma)=81$,
$k=6000$ chains, 4000 steps. Agreement is within $0.3\%$ over a factor of twenty in
$\varepsilon$.}
\label{tab:bias}
\end{table}

\begin{table}[t]
\centering
\begin{tabular}{rrr}
\toprule
$\varepsilon$ & $R$ uncorrected & $R$ corrected \\
\midrule
0.20 & 1.1087 & 0.9978 \\
0.10 & 1.0495 & 0.9970 \\
0.05 & 1.0236 & 0.9980 \\
0.01 & 1.0042 & 0.9992 \\
\bottomrule
\end{tabular}
\caption{Effect of the correction \eqref{eq:correction}. Besides removing the bias, it makes the
result largely insensitive to the choice of $\varepsilon$, which removes the main tuning burden.
Enabled by default (\code{bias\_correct=True}).}
\label{tab:corr}
\end{table}

\section{Why the SVGD repulsion fails: Gram-matrix degeneracy}
\label{sec:gram}

SVGD's spread comes from the repulsion term $\sum_j \nabla_{x_j} k(x_j, x_i)$, whose magnitude is
controlled by the off-diagonal entries of the Gram matrix. With the RBF kernel
$k(x,x') = \exp(-\|x-x'\|^2/h)$ and the usual median heuristic
$h = \operatorname{median}_{i<j}\|x_i-x_j\|^2 / \log k$, high dimension is unkind. Squared
distances concentrate---for a roughly isotropic ensemble $\|x_i-x_j\|^2$ has relative standard
deviation $\mathcal{O}(d^{-1/2})$---so for essentially every pair
$\|x_i-x_j\|^2/h \approx \log k$, giving
\begin{equation}
  k(x_i,x_j) \;\approx\; e^{-\log k} \;=\; \frac{1}{k}\quad (i\neq j),
  \qquad k(x_i,x_i) = 1 .
\end{equation}
The Gram matrix is thus $K \approx I + \tfrac1k(\mathbf{1}\mathbf{1}^{\!\top}-I)$. Its row-wise
participation ratio, the effective number of neighbours each particle interacts with, is
\begin{equation}
  n_{\text{eff}} \;=\; \frac{\big(\sum_j K_{ij}\big)^2}{\sum_j K_{ij}^2}
  \;\longrightarrow\; \frac{2^2}{1} \;=\; 4 \qquad (k \to \infty),
\end{equation}
independently of $k$. Each particle effectively sees about four neighbours no matter how large the
ensemble, the repulsion is negligible, the drift dominates, and the ensemble collapses onto the
mode. Table~\ref{tab:gram} confirms all of this numerically, including the prediction that adding
particles does not help.

An inverse multiquadric kernel $(1+\|x-x'\|^2/h)^{-1/2}$ decays only as $(1+\log k)^{-1/2}$ and so
does not degenerate; it is an effective alternative fix, but it must still be run in whitened
coordinates and it inherits SVGD's contraction behaviour. Langevin avoids the issue entirely by
not using a kernel.

\begin{table}[t]
\centering
\begin{tabular}{rrrrrrr}
\toprule
$k$ & mean $\|x_i-x_j\|^2/h$ & $\log k$ & RBF off-diag. & $1/k$ & IMQ off-diag. & $n_{\text{eff}}$ \\
\midrule
 200 & 5.310 & 5.298 & $5.34\times10^{-3}$ & $5.0\times10^{-3}$ & 0.399 & 4.26 \\
 800 & 6.696 & 6.685 & $1.40\times10^{-3}$ & $1.25\times10^{-3}$ & 0.361 & 4.56 \\
3200 & 8.087 & 8.071 & $3.73\times10^{-4}$ & $3.1\times10^{-4}$ & 0.332 & 4.93 \\
\bottomrule
\end{tabular}
\caption{Measured Gram-matrix statistics on whitened ensembles, $d=325$. The predictions
$\|x_i-x_j\|^2/h \to \log k$, RBF off-diagonal $\to 1/k$, and $n_{\text{eff}} \to 4$ all hold. The
effective neighbour count does not grow with $k$, so increasing the ensemble size cannot repair
the repulsion.}
\label{tab:gram}
\end{table}

\section{The algorithm}

\begin{algorithm}[H]
\SetAlgoLined
\KwIn{log-density $\log p$, start $x_0$, chains $k$, steps $T$, warmup $T_{\text{warm}}$, step size $\varepsilon$}
\KwOut{$\{x_i\}_{i=1}^k$ approximately distributed as $p$}
$\xmap \leftarrow \arg\max \log p$ \tcp*{mode-find from $x_0$}
$H \leftarrow -\nabla^2 \log p(\xmap)$; eigendecompose $H = V\Lambda V^{\!\top}$ in float64\;
\lIf{$\min \operatorname{diag}(\Lambda) \le 0$}{\textbf{abort}: metric undefined (see \S\ref{sec:failure})}
$L \leftarrow V \Lambda^{-1/2} V^{\!\top}$ \tcp*{$= H^{-1/2}$, symmetric}
$y_i \sim \mathcal{N}(0, I_d)$ for $i=1,\dots,k$ \tcp*{$=$ the Laplace approximation}
$(\ell_i, g_i) \leftarrow$ value-and-grad of $\log p$ at $\xmap + y_i L$; \ \ $g_i \leftarrow g_i L$\;
\For{$t = 1$ \KwTo $T_{\text{warm}} + T$}{
  $z_i \sim \mathcal{N}(0, I_d)$; \ \ $s_i \leftarrow \varepsilon g_i + \sqrt{2\varepsilon}\, z_i$\;
  $(\ell_i', g_i') \leftarrow$ value-and-grad at $\xmap + (y_i + s_i)L$; \ \ $g_i' \leftarrow g_i' L$\;
  \eIf{\normalfont metropolis}{
    $\Lambda_i \leftarrow (\ell_i' - \ell_i) + \tfrac12\|z_i\|^2 - \|s_i + \varepsilon g_i'\|^2/(4\varepsilon)$\;
    accept chain $i$ with prob.\ $\min(1, e^{\Lambda_i})$, updating $(y_i,\ell_i,g_i)$; tally per chain\;
    \lIf{$t \le T_{\text{warm}}$}{$\log\varepsilon \mathrel{+}= \gamma(\hat a - 0.574)$ \tcp*[f]{else freeze}}
  }{
    $y_i \leftarrow y_i + s_i$; \ \ \lIf{$y$ non-finite}{\textbf{abort}: transient (see \S\ref{sec:nongauss})}
  }
}
\lIf{\normalfont metropolis}{\textbf{abort} if $\Lambda_F > 5\%$ or $F > 10\%$ (see \S\ref{sec:sticky})}
\lElse{$y_i \leftarrow \bar y + (y_i - \bar y)\sqrt{1-\varepsilon/2}$ \tcp*[f]{\eqref{eq:correction}}}
\Return $\xmap + y_i L$\;
\caption{Whitened Langevin, Metropolis-adjusted by default (\code{MSVGD.whitened\_ula})}
\end{algorithm}

\paragraph{On the initialization.} Drawing $y \sim \mathcal{N}(0,I)$ is exactly the Laplace
approximation expressed in whitened coordinates, so the chains begin at approximate stationarity
and there is no burn-in to discard---which is why a couple of thousand steps suffices. Langevin is
ergodic and does forget its initialization: starting instead from a tight ball reaches an
identical answer by 2000 steps, merely more slowly at first (Table~\ref{tab:init}). What is
\emph{not} safe is starting wider, since the initial ensemble then already straddles the region
where the quadratic model fails. $\mathcal{N}(0,I)$ is both the natural choice and the safe upper
end.

\begin{table}[t]
\centering
\begin{tabular}{lrrr}
\toprule
initialization & 50 steps & 200 steps & 2000 steps \\
\midrule
$\mathcal{N}(0,I)$ \ (Laplace) & 0.596 & 0.100 & \textbf{0.0484} \\
$0.05\cdot\mathcal{N}(0,I)$ \ (tight ball) & 0.983 & 0.109 & \textbf{0.0484} \\
$3\cdot\mathcal{N}(0,I)$ \ (over-dispersed) & \multicolumn{3}{c}{diverged} \\
\bottomrule
\end{tabular}
\caption{Energy distance to the reference (sampling floor $0.048$) versus initialization and
number of steps, $d=325$. The Laplace start converges faster but not to a different answer.}
\label{tab:init}
\end{table}

\section{Empirical validation}

Table~\ref{tab:main} compares the method against an 8-chain, $64{,}000$-draw NUTS gold standard
(0 divergences, $\hat R \le 1.0048$, minimum ESS 1202) on the $d=325$ FitzHugh--Nagumo posterior.
Metrics are: energy distance to the reference in reference-whitened coordinates; the
variance-weighted ratio of ensemble to reference variance along the reference principal axes,
which unlike energy distance is sensitive to the few directions carrying most of the variance;
the Stein diagnostic \eqref{eq:steinR}; the Mahalanobis error of the ensemble mean; and the mean
absolute deviation of the three parameter credible-interval widths. Every ``floor'' is measured,
by scoring an independent $800$-draw subsample of the reference through the identical pipeline.

\begin{table}[t]
\centering
\small
\begin{tabular}{lrrrrrr}
\toprule
& energy & varwtd & $R$ & mean err. & CI dev.\ (\%) & wall-clock \\
\midrule
NUTS, $k=800$ \emph{(sampling floor)} & 0.048 & 1.020 & 1.003 & 0.034 & 3.5 & 125\,s \\
\textbf{whitened ULA} & \textbf{0.050} & 0.972 & 0.993 & 0.037 & \textbf{2.3} & \textbf{3.6\,s} \\
SVGD, density-reweighted kernel & 3.665 & 1.409 & 0.244 & 0.176 & 5.0 & $\sim$4\,s \\
SVGD, standard kernel & 6.742 & 0.674 & 0.022 & 0.068 & 27.4 & $\sim$4\,s \\
\bottomrule
\end{tabular}
\normalsize
\caption{Whitened ULA against the NUTS gold standard and against the two SVGD configurations it
replaces. The method is at the sampling floor on every metric. Note that the credible-interval
deviation column orders the SVGD rows favourably while their actual distributional error is
$70$--$140\times$ worse, which is why several complementary metrics are reported.}
\label{tab:main}
\end{table}

A caution on the reweighted-kernel row: its per-coordinate standard deviations average to $1.037$,
which looks nearly correct, but this is cancellation. It is $43$--$56\%$ \emph{over}-dispersed in
the leading eigendirections and $88\%$ \emph{collapsed} in the remaining 305. Single scalar
summaries conceal this; the variance-weighted column does not.

\section{Robustness to non-Gaussianity}
\label{sec:nongauss}

Because $H^{-1}$ enters only as a preconditioner and never as the answer, ULA targets the exact
posterior whatever its shape. Non-Gaussianity therefore does not \emph{bias} the method the way it
breaks any procedure that returns a Gaussian. Three things can still go wrong, and they differ
sharply in severity.

\paragraph{(a) Bias --- does not occur.} On a $12$-dimensional banana
$\log p \propto -x_0^2/2S^2 - \tfrac12\big(x_1 - b(x_0^2-S^2)\big)^2 - \tfrac12\|x_{2:}\|^2$
($S=10$, $b=0.05$), the Laplace approximation is badly wrong: it assigns $x_1$ a standard
deviation $0.14$ times the truth. Whitened ULA nonetheless walks toward the correct answer,
the skewness of $x_1$ climbing $0.75 \to 1.43 \to 2.00$ against an exact value of $2.79$ as the
budget grows. It is slow, not wrong.

\paragraph{(b) Mixing --- degrades gracefully.} The speedup of
Section~\ref{sec:analysis} assumes $H$ at the mode represents the global geometry. Where it does
not---curved ridges, multimodality---the preconditioner is merely unhelpful rather than harmful,
and the sampler reverts toward the unwhitened mixing rate.

\paragraph{(c) Transience --- the real hazard, and it is specific to polynomial ODEs.}
The explicit Euler step of \eqref{eq:ula} is unstable for super-quadratic potentials
\cite{robertstweedie}. For $U(x) = x^2/2 + a|x|^p/p$ the noise-free map is
$x \mapsto x(1-\varepsilon) - \varepsilon a\,x^{p-1}$, which is expanding once
$\varepsilon\big(1 + a x^{p-2}\big) > 2$, i.e.\ beyond a critical radius
\begin{equation}
  R_{\text{crit}} \;=\; \Big(\frac{2/\varepsilon - 1}{a}\Big)^{\!1/(p-2)}
  \;\propto\; \varepsilon^{-1/(p-2)} .
  \label{eq:rcrit}
\end{equation}
A chain that wanders past $R_{\text{crit}}$ is flung outward and never returns; the algorithm
diverges rather than returning a biased answer. Table~\ref{tab:transience} shows divergence
occurring precisely when $R_{\text{crit}}$ falls inside the range the chains actually visit.

The scaling in \eqref{eq:rcrit} is the practical problem. Reducing the step size buys safety only
as $\varepsilon^{-1/(p-2)}$: at $p=4$ that is $\varepsilon^{-1/2}$, but at $p=6$ it is
$\varepsilon^{-1/4}$, so cutting $\varepsilon$ by a factor of $20$ enlarges the safe region by
only $2.1\times$. \textbf{At high polynomial degree the step size is a very weak lever.}

\begin{table}[t]
\centering
\begin{tabular}{rrrlr}
\toprule
$p$ & $a$ & $\varepsilon$ & outcome & $R_{\text{crit}}$ \\
\midrule
2 & $0$        & any   & finite, s.d.\ $1.00$        & $\infty$ \\
4 & $10^{-3}$  & any   & finite, s.d.\ $1.00$        & 95--446 \\
4 & $1$        & 0.20  & \textbf{diverged}           & 3.00 \\
4 & $1$        & 0.05  & finite, $\max|x| = 2.32$    & 6.24 \\
4 & $1$        & 0.01  & finite, $\max|x| = 2.14$    & 14.11 \\
6 & $10^{-3}$  & any   & finite, s.d.\ $1.00$        & 9.7--21 \\
6 & $1$        & \emph{all three} & \textbf{diverged} & 1.73--3.76 \\
\bottomrule
\end{tabular}
\caption{Transience of ULA on $U = x^2/2 + a|x|^p/p$ in $d=8$; the potential is Gaussian at the
mode, so the positive-definiteness check passes and the failure is purely a discretization
effect. Divergence occurs exactly when $R_{\text{crit}}$ from \eqref{eq:rcrit} enters the range
the chains visit ($\approx 2$--$4.5$).}
\label{tab:transience}
\end{table}

\subsection*{Why this makes ODE posteriors adversarial}

The MAGI log-density contains the term $r^{\!\top} K^{-1} r$ with
$r = f(X,\theta) - \dot\mu - m(X-\mu)$. For a vector field $f$ polynomial of degree $q$ in the
state, $r$ is degree $q$ and this term is \textbf{degree $2q$}. FitzHugh--Nagumo has a cubic
term $V^3/3$, so $q=3$ and the potential is degree six---the worst row of
Table~\ref{tab:transience}.

Measuring the growth of $U$ along whitened directions confirms this, and shows where it lives:

\begin{center}
\begin{tabular}{lrrrr}
\toprule
direction & $U(3)$ & $U(30)$ & $U(100)$ & fitted exponent \\
\midrule
flattest eigenvector & $5.13\times10^{1}$ & $6.90\times10^{5}$ & $3.66\times10^{9}$ & \textbf{6.20} \\
median eigenvector   & $3.74\times10^{0}$ & $4.41\times10^{2}$ & $5.06\times10^{3}$ & 2.03 \\
stiffest eigenvector & $4.06\times10^{0}$ & $4.46\times10^{2}$ & $4.99\times10^{3}$ & 2.01 \\
\bottomrule
\end{tabular}
\end{center}

The potential is quadratic in almost every direction---the Gaussian-process prior and the
observation term dominate there---but degree six along the \emph{flattest} eigenvector, which is
precisely the direction the whitening stretches most, by $\lambda_{\min}^{-1/2}$.

This sharpens the account of the $\sigma = 0.5$ failure in Section~\ref{sec:failure}. It is not
merely that the metric is inaccurate. Weak identification shrinks $\lambda_{\min}$, whitening
therefore stretches the flat direction by a large factor, the chains reach the region where the
degree-six term dominates, and there the Euler discretization is transient. The pre-flight
diagnostic works because $\lambda_{\min}$ controls how far the chains are thrown along exactly
the direction in which the potential is worst behaved.

\paragraph{Consequences.} FitzHugh--Nagumo at the baseline data density is safe because the
observation term is strong enough to keep the chains near the data, far inside $R_{\text{crit}}$.
It is, however, already at $p=6$: a quintic vector field would give $p=10$, where
$R_{\text{crit}} \propto \varepsilon^{-1/8}$ and the step size ceases to be a usable control at
all. Shrinking $\varepsilon$ is therefore not the remedy; the Metropolis correction of
Section~\ref{sec:mala} is, but only partially, as follows.

\subsection{The Metropolis correction cures transience but introduces stickiness}
\label{sec:sticky}

MALA rejects the overshooting proposal, so no chain is thrown to infinity. It does \emph{not}
make the sampler reliable on these targets, because MALA is not geometrically ergodic when the
tails are lighter than exponential \cite{robertstweedie}: from far out, \emph{every} proposal
overshoots and is rejected, so a chain that \emph{starts} there can never move. Transience is
replaced by its mirror image---frozen chains---and where transience announced itself with a NaN,
stickiness returns a plausible-looking but over-dispersed answer.

The effect is entirely one of initialization, as the following isolates. On
$U = x^2/2 + |x|^p/p$ in $d=8$ with $k=4000$:

\begin{center}
\begin{tabular}{lrrr}
\toprule
$p$ & frozen chains & $\E[x^2]$ / exact & mean acceptance \\
\midrule
2  & $0.00\%$  & 0.993 & 0.571 \\
4  & $0.03\%$  & 0.996 & 0.573 \\
6  & $16.35\%$ & 1.50  & 0.577 \\
10 & $35.68\%$ & 2.57  & 0.574 \\
\midrule
\multicolumn{4}{l}{\emph{$p=10$, but initialized at $0.3\cdot\mathcal{N}(0,I)$, inside the bulk:}} \\
10 & --- & \textbf{0.9989} & --- \\
\bottomrule
\end{tabular}
\end{center}

Started inside the bulk, MALA recovers $\E[x^2]$ to $0.1\%$ even at $p=10$; started at
$\mathcal{N}(0,I)$, a third of the chains freeze in the tail and contribute $63\%$ of the second
moment. Note the mean acceptance rate: it sits at the $0.574$ target throughout and is
completely uninformative about the failure. This is why the implementation tracks acceptance
\emph{per chain}.

\paragraph{Detection.} Two independent things can go wrong, so the implementation raises on
either. Writing $F$ for the fraction of chains that accepted nothing at all and $\Lambda_F$ for
the relative change in the ensemble's scale from dropping them:
\begin{itemize}[leftmargin=1.3em,itemsep=2pt]
  \item $\Lambda_F > 5\%$: the frozen chains are tail outliers dominating the moments.
  \item $F > 10\%$: that much of the ``sample'' is just the initialization. This is the case a
  leverage test cannot see, because when the \emph{whole} ensemble starts outside the bulk the
  frozen chains are typical of it and $\Lambda_F$ stays small.
\end{itemize}
Table~\ref{tab:guard} shows the two triggers separating every case tested, with no false alarm
on the well-identified problem.

\begin{table}[t]
\centering
\begin{tabular}{lrrll}
\toprule
case & $F$ & $\Lambda_F$ & verdict & accuracy \\
\midrule
synthetic $p=2$    & $0.00\%$  & $0.00\%$  & passed & $\E[x^2]$ ratio 0.9933 \\
synthetic $p=4$    & $0.03\%$  & $0.10\%$  & passed & $\E[x^2]$ ratio 0.9964 \\
synthetic $p=6$    & $16.35\%$ & $17.76\%$ & \textbf{raised} & --- \\
synthetic $p=10$   & $35.68\%$ & $25.58\%$ & \textbf{raised} & --- \\
MAGI baseline, 3 seeds & $0.5$--$1.3\%$ & $0.1$--$0.3\%$ & passed & energy $0.0496$--$0.0503$ \\
MAGI $\sigma=0.5$  & $42.58\%$ & $0.48\%$  & \textbf{raised} & --- \\
\bottomrule
\end{tabular}
\caption{The two triggers, across every case tested. The $\sigma=0.5$ row is the one that
motivates the second trigger: $43\%$ of chains frozen but negligible leverage, because there the
entire ensemble is stuck near its start.}
\label{tab:guard}
\end{table}

\paragraph{What actually fixes it.} Nothing here does, for a genuinely weakly identified target;
both variants fail and both are detected. A tamed or projected integrator that caps the drift
would remove the transience at its source, and a better-centred initialization would remove the
stickiness, but the underlying problem is that $H^{-1}$ is the wrong metric in that regime. The
honest scope is the pre-flight check of Section~\ref{sec:failure}.

\section{Failure mode and diagnostics}
\label{sec:failure}

The metric is only as trustworthy as the quadratic model at $\xmap$. If the posterior has a
near-flat direction, $H^{-1}$ claims very large variance there, whitening stretches that direction
accordingly, and the chains are placed where the model's nonlinearity---here the cubic term of the
FitzHugh--Nagumo vector field---no longer resembles a quadratic. The method then diverges.

This is predictable from the spectrum of $H$ alone, at no additional cost and with no reference
chain (Table~\ref{tab:failure}). Both failure modes raise rather than return silently:

\begin{itemize}[leftmargin=1.3em,itemsep=2pt]
  \item \textbf{$H$ not positive definite.} Caught before any sampling. In practice this most
  often signals an unconverged mode-find, so the check doubles as a MAP-convergence test.
  \item \textbf{Non-finite state during sampling.} Caught during; the raised error reports
  $\lambda_{\min}$ and $\tr(H^{-1})$ and identifies weak identification as the cause.
\end{itemize}

\begin{table}[t]
\centering
\begin{tabular}{lrrrrl}
\toprule
setting & $\lambda_{\min}$ & $\operatorname{cond}(H)$ & $\tr(H^{-1})$ & flattest s.d. & outcome \\
\midrule
baseline, $\sigma=0.2$ & 5.307 & $1.9\times10^{4}$ & 0.98 & 0.43 & at the floor \\
half the observations & 1.784 & $7.4\times10^{4}$ & 1.68 & 0.75 & works \\
quarter of the observations & 1.591 & $1.0\times10^{5}$ & 1.88 & 0.79 & degraded \\
noisy, $\sigma=0.5$ & \textbf{0.0078} & $3.1\times10^{7}$ & \textbf{133.2} & 11.36 & \textbf{diverges} \\
\bottomrule
\end{tabular}
\caption{Pre-flight diagnostic. As identification weakens, $\lambda_{\min}$ collapses and
$\tr(H^{-1})$---the total variance claimed by the Laplace approximation---explodes. At
$\sigma=0.5$ the quadratic model predicts a standard deviation of $11.4$ in the flattest
direction, where the true posterior is bounded by the ODE nonlinearity rather than by curvature.}
\label{tab:failure}
\end{table}

Where the check fails, NUTS on the untransformed target remains the appropriate fallback; nothing
described here improves on it in that regime. It is worth noting, however, that this method fails
\emph{loudly}, whereas the SVGD configurations at the same setting return plausible-looking but
badly wrong ensembles.

\section{Implementation notes}

\begin{itemize}[leftmargin=1.3em,itemsep=3pt]
  \item \textbf{The eigendecomposition is performed in float64} on the host irrespective of the
  particle dtype. Whitening is governed by the \emph{smallest} eigenvalues of $H$, since
  $H^{-1/2}$ amplifies the flattest directions most, and those are precisely the ones a float32
  decomposition resolves worst.

  \item \textbf{The metric admits almost no error.} Replacing $H$ by
  $0.99H + 0.01\operatorname{diag}(H)$ costs an order of magnitude in accuracy; evaluating $H$ one
  posterior standard deviation away from $\xmap$---a relative Frobenius perturbation of
  $0.4\%$---is enough to make it indefinite. Diagonal, block-diagonal and prior-only
  approximations all fail, because the ill-conditioning resides in the correlations rather than in
  the marginal scales ($\operatorname{cond}(\operatorname{diag} H)=30$ against
  $\operatorname{cond}(H)=1.9\times10^4$).

  \item \textbf{Cost} is $\mathcal{O}(d)$ gradient evaluations, $\mathcal{O}(d^2)$ memory and
  $\mathcal{O}(d^3)$ to factor, which is what bounds the accessible dimension---comfortable to
  $d\sim 10^3$--$10^4$. Scaling beyond a dense Hessian is the principal open problem; the natural
  randomized low-rank route is unpromising precisely because the method depends on the bottom of
  the spectrum.

  \item \textbf{The sampling loop} is a single \code{lax.scan}, so the whole run compiles to one
  program. Setup dominates: $2.1$\,s for the mode-find, $0.8$\,s for the Hessian, and under
  $0.5$\,s for 2000 Langevin steps at $k=800$, $d=325$. The Metropolis correction is
  approximately free, since \code{value\_and\_grad} costs what the gradient alone costs and the
  accepted state is carried between steps.

  \item \textbf{Acceptance is tracked per chain, not just on average.} The characteristic
  failure of Section~\ref{sec:sticky} leaves the mean acceptance sitting exactly at its target
  while a third of the chains have never moved, so a mean-acceptance check---the obvious
  thing---detects nothing.
\end{itemize}

\begin{thebibliography}{1}
\bibitem{robertsrosenthal} G.~O. Roberts and J.~S. Rosenthal.
Optimal scaling of discrete approximations to Langevin diffusions.
\emph{J. R. Statist. Soc. B}, 60(1):255--268, 1998.
\bibitem{robertstweedie} G.~O. Roberts and R.~L. Tweedie.
Exponential convergence of Langevin distributions and their discrete approximations.
\emph{Bernoulli}, 2(4):341--363, 1996.
\end{thebibliography}

\end{document}
